% ---------------------------------------------------------------------------
% Chapitre 01 — Introduction : pourquoi analyser statiquement React ?
% Sources : notes/16 (§1, §6), notes/10 et 11 (bug React visé), notes/13 (§6.1),
% notes/15 (§5.1), README.md, docs/usage.md. Sorties observées au commit e67b10a
% (binaire target/debug/reactant), fichiers d'exemple sous /tmp/ch01/.
% ---------------------------------------------------------------------------
\chapter{Introduction : pourquoi analyser statiquement React ?}
\label{chap:introduction}

\begin{objectifs}
  \item Reconnaître trois défauts React typiques et ce qui se passe réellement à l'exécution.
  \item Comprendre pourquoi un linter syntaxique ne suffit pas.
  \item Saisir l'idée directrice : interpréter abstraitement un composant.
  \item Connaître la promesse de l'outil : trois niveaux de diagnostic et la soundness.
  \item Lancer \reactant{} et lire ses sorties humaine et JSON.
  \item Situer le projet (taille, règles, décisions) et choisir un parcours de lecture.
\end{objectifs}

\prerequis{aucun. Il suffit de savoir écrire un composant React avec \code{useState} et \code{useEffect}, et de lire du Rust.}

Un composant React est une fonction courte. On la lit en quelques secondes, on
la teste en quelques minutes, et pourtant les défauts qu'elle contient sont
parmi les plus coûteux d'une application web : un onglet qui gèle, un écran qui
affiche les données d'un autre utilisateur, un chronomètre qui reste bloqué.
Ces défauts ont un point commun : ils ne se voient pas dans le texte du
composant, mais dans la \emph{suite de ses exécutions}. Ce chapitre explique
pourquoi il faut, pour les trouver, raisonner sur toutes ces exécutions à la
fois, et comment \reactant{} s'y prend. Il reste volontairement informel : les
chapitres suivants reprennent chaque idée avec sa définition exacte et le code
qui l'incarne.

\section{Trois programmes React fautifs}
\label{sec:introduction-trois-bugs}

Commençons par trois composants très simples. Chacun compile, s'affiche, et
passe un test qui vérifie le premier rendu. Chacun est pourtant faux. Pour
comprendre en quoi, il faut d'abord rappeler dans quel ordre React exécute un
composant.

\begin{react}[le cycle rendu, commit, effets]
  React exécute un composant en trois temps~\cite{ReactRenderCommit}.
  \begin{enumerate}
    \item \textbf{Déclenchement.} Un rendu est demandé soit parce que le
      composant apparaît pour la première fois (le \emph{montage}), soit parce
      qu'un état a été mis à jour par un setter (\code{setN(...)}) dans ce
      composant ou dans l'un de ses ancêtres.
    \item \textbf{Rendu.} React \emph{appelle} la fonction du composant. Elle
      lit l'état courant, calcule le JSX et ne doit rien faire d'autre : le
      rendu est un calcul pur.
    \item \textbf{Commit.} React reporte dans le DOM les différences entre
      l'ancien et le nouveau JSX.
  \end{enumerate}
  Viennent ensuite les \textbf{effets} : après le commit, React exécute les
  fonctions passées à \code{useEffect}. Sans tableau de deps, un effet
  s'exécute après \emph{chaque} commit ; avec \code{[]}, une seule fois après
  le montage ; avec \code{[a, b]}, seulement si l'un des deps a changé depuis
  le rendu précédent, la comparaison se faisant par \code{Object.is}. Enfin,
  un setter ne modifie pas la variable du rendu en cours : il \emph{planifie}
  un nouveau rendu, dans lequel \code{useState} renverra la nouvelle valeur.
  Si cette valeur est identique à l'ancienne (encore au sens de
  \code{Object.is}), React peut s'épargner le re-rendu.
\end{react}

\Cref{fig:introduction-cycle-react} résume ce cycle. La flèche qui remonte des
effets vers le rendu est la source de tous les ennuis de ce chapitre : un effet
qui appelle un setter relance le cycle.

\begin{figure}[htbp]\centering
\begin{tikzpicture}[node distance=9mm and 11mm]
  \node[bloc] (trig) {déclenchement\\{\scriptsize montage ou setter}};
  \node[bloc,right=of trig] (render) {rendu\\{\scriptsize appel du composant}};
  \node[bloc,right=of render] (commit) {commit\\{\scriptsize mise à jour du DOM}};
  \node[bloc,right=of commit] (eff) {effets\\{\scriptsize \texttt{useEffect}}};
  \draw[fleche] (trig) -- (render);
  \draw[fleche] (render) -- (commit);
  \draw[fleche] (commit) -- (eff);
  \draw[flechemay] (eff.south) -- ++(0,-8mm) -| node[etiquette,pos=0.25,below]{un effet appelle un setter : nouveau rendu planifié} (render.south);
  \draw[flechemay] (eff.north) -- ++(0,7mm) -| node[etiquette,pos=0.25,above]{un handler (clic, frappe\dots) appelle un setter} (trig.north);
\end{tikzpicture}
\caption{Le cycle d'exécution d'un composant React. Les flèches pointillées
sont les retours possibles : un effet ou un handler qui appelle un setter
planifie un nouveau rendu.}
\label{fig:introduction-cycle-react}
\end{figure}

\subsection{Premier défaut : la boucle infinie}

\begin{lstlisting}[language=TypeScript,caption={\code{Counter} : un effet sans deps qui écrit l'état qu'il lit (\file{/tmp/ch01/loop.tsx})},label={lst:introduction-loop}]
import { useEffect, useState } from "react";

export function Counter() {
  const [n, setN] = useState(0);
  useEffect(() => {
    setN(n + 1);
  });
  return <p>{n}</p>;
}
\end{lstlisting}

Déroulons la \cref{lst:introduction-loop} à la main. Au montage, \code{n} vaut
\code{0} ; le rendu produit \code{<p>0</p>}, React le commit, puis exécute
l'effet, qui appelle \code{setN(1)}. Un rendu est planifié : \code{n} vaut
\code{1}, commit, et comme l'effet n'a pas de tableau de deps, il s'exécute de
nouveau et appelle \code{setN(2)}. Et ainsi de suite : \code{n} prend les
valeurs $0, 1, 2, 3, \dots$ sans fin, chaque valeur étant différente de la
précédente, si bien que React ne peut jamais s'épargner le re-rendu. Le
navigateur passe son temps à rendre ce composant ; en développement, React
finit par se plaindre dans la console d'un nombre excessif de mises à jour
imbriquées.

Ce défaut est \emph{certain} : aucune entrée, aucun clic, aucune prop ne
l'empêche. Il n'apparaît pourtant pas dans un test qui ne regarde que le premier
rendu, puisque celui-ci est correct.

\subsection{Deuxième défaut : la dépendance oubliée}

\begin{lstlisting}[language=TypeScript,caption={\code{Profile} : l'effet lit \code{userId} mais ne le déclare pas (\file{/tmp/ch01/deps.tsx})},label={lst:introduction-deps}]
import { useEffect, useState } from "react";

export function Profile({ userId }: { userId: string }) {
  const [user, setUser] = useState(null);
  useEffect(() => {
    fetch(`/api/users/${userId}`)
      .then((r) => r.json())
      .then(setUser);
  }, []);
  return <p>{user ? user.name : "..."}</p>;
}
\end{lstlisting}

Ici, pas de boucle. Au montage, l'effet charge l'utilisateur \code{userId} et
l'affiche. Mais le tableau de deps est vide : l'effet ne s'exécutera \emph{plus
jamais}. Si le parent rend ensuite \code{<Profile userId="bob" />} à la place
de \code{<Profile userId="alice" />}, le composant est re-rendu avec la
nouvelle prop, mais l'effet ne tourne pas, et l'écran continue d'afficher
Alice. Rien ne plante ; l'interface ment, en silence.

Ce défaut-là n'est pas certain au même sens que le premier : il ne se manifeste
que si \code{userId} change effectivement pendant la vie du composant, ce que le
composant seul ne permet pas de savoir.

\subsection{Troisième défaut : la fermeture figée}

\begin{lstlisting}[language=TypeScript,caption={\code{Timer} : le callback de \code{setInterval} lit une valeur figée au montage (\file{/tmp/ch01/timer.tsx})},label={lst:introduction-timer}]
import { useEffect, useState } from "react";

export function Timer() {
  const [seconds, setSeconds] = useState(0);
  useEffect(() => {
    const id = setInterval(() => {
      setSeconds(seconds + 1);
    }, 1000);
    return () => clearInterval(id);
  }, []);
  return <p>{seconds} s</p>;
}
\end{lstlisting}

L'effet de la \cref{lst:introduction-timer} s'exécute une fois, au montage, et
installe un intervalle. Le callback passé à \code{setInterval} est une
fermeture (\emph{closure}) : il capture la variable \code{seconds}
\emph{du rendu de montage}, qui vaut \code{0}, et la gardera toute sa vie. À la
première seconde, il appelle \code{setSeconds(0 + 1)} : l'écran affiche
\og 1 s\fg{}. À la deuxième, il appelle encore \code{setSeconds(0 + 1)}, car
sa variable \code{seconds} vaut toujours \code{0}. L'état reste à \code{1},
React constate par \code{Object.is} que rien n'a changé, et le chronomètre
reste bloqué à \og 1 s\fg{} pour toujours. On parle de \terme{stale closure}
(fermeture périmée).

La correction consiste à ne plus lire la variable capturée, mais à passer au
setter une \emph{fonction de mise à jour} (un \emph{updater}) :
\code{setSeconds((s) => s + 1)}. React appelle cette fonction avec la valeur
courante de l'état, et non celle du montage.

\subsection{Ce que ces trois exemples ont en commun}

\Cref{tab:introduction-trois-bugs} met les trois défauts côte à côte. Aucun
n'est visible au premier rendu ; tous dépendent de la \emph{succession} des
rendus, commits et effets ; et chacun demande de savoir quelque chose sur
\emph{toutes} les exécutions possibles du composant : que \code{n} ne se
stabilise jamais, que \code{userId} peut changer, que \code{seconds} lu par le
callback est toujours la valeur du montage.

\begin{table}[htbp]\centering\small
\begin{tabularx}{\linewidth}{@{}lXXl@{}}
  \toprule
  Composant & Symptôme à l'exécution & Condition d'apparition & Nature \\
  \midrule
  \code{Counter} & rendus sans fin, onglet saturé & toujours, dès le montage & certain \\
  \code{Profile} & données d'un autre utilisateur & si la prop \code{userId} change & possible \\
  \code{Timer}   & affichage bloqué à \og 1 s\fg{} & toujours, dès la deuxième seconde & certain \\
  \bottomrule
\end{tabularx}
\caption{Les trois défauts d'ouverture. La dernière colonne annonce la
distinction entre défaut certain et défaut possible, qui structure les niveaux
de diagnostic (\cref{sec:introduction-promesse}).}
\label{tab:introduction-trois-bugs}
\end{table}

\section{Ce qu'un linter syntaxique voit, et ce qu'il ne voit pas}
\label{sec:introduction-linter}

L'outil que tout projet React installe d'abord est un \emph{linter} :
\code{eslint-plugin-react-hooks}, avec ses deux règles historiques
\code{rules-of-hooks} et \code{exhaustive-deps}~\cite{ReactRulesOfHooks}. Il
est rapide, il tourne dans l'éditeur, et il attrape une partie des défauts
précédents. Pour comprendre ce qu'il ne peut pas attraper, il faut préciser ce
qu'il regarde.

\begin{definition}{Analyse syntaxique d'un composant}{introduction-syntaxique}
  Une \terme{analyse syntaxique} (ou \emph{lexicale}) décide à partir de la
  \emph{forme} du texte : elle parcourt l'arbre de syntaxe d'une fonction, dans
  un fichier, et reconnaît des motifs (un appel \code{useX} sous un \code{if},
  un identifiant lu dans un effet et absent du tableau de deps). Elle ne
  calcule pas les \emph{valeurs} que prennent les variables quand le programme
  s'exécute.
\end{definition}

\subsection{Ce que la forme suffit à voir}

Deux des défauts de la \cref{sec:introduction-trois-bugs} sont écrits \og en
toutes lettres\fg{}. Dans \code{Profile} (\cref{lst:introduction-deps}),
l'identifiant \code{userId} apparaît dans le corps de l'effet et pas dans
\code{[]} : il suffit de comparer deux listes de noms, et c'est exactement ce
que fait \code{exhaustive-deps}. De même, un hook appelé après un retour
anticipé est reconnaissable sur l'arbre de syntaxe :

\begin{lstlisting}[language=TypeScript,caption={\code{Panel} : un \code{useState} après un retour anticipé (\file{/tmp/ch01/cond.tsx})},label={lst:introduction-cond}]
import { useState } from "react";

export function Panel({ open }: { open: boolean }) {
  if (!open) return null;
  const [tab, setTab] = useState(0);
  return <button onClick={() => setTab(tab + 1)}>{tab}</button>;
}
\end{lstlisting}

\begin{react}[les hooks sont identifiés par leur rang]
  React ne connaît pas le \emph{nom} des variables d'état : il range les hooks
  d'un composant dans une liste, dans l'ordre où ils sont appelés, et les
  retrouve au rendu suivant par leur \emph{position}~\cite{ReactRulesOfHooks}.
  Si un rendu appelle un hook de moins (ici, quand \code{open} vaut
  \code{false}), tous les hooks suivants sont décalés d'un cran. D'où la règle :
  les hooks s'appellent au premier niveau du composant, jamais sous une
  condition, une boucle ou après un \code{return} conditionnel.
\end{react}

\reactant{} signale ces deux composants, comme le ferait le linter :

\begin{sortie}
  Panel  (2 hooks)  /tmp/ch01/cond.tsx
    error  conditional-hook  [hook:0]  (line 5:8)  this hook is called conditionally (not on every render path)
       (1 trace step(s), rerun with --trace)

  Profile  (2 hooks)  /tmp/ch01/deps.tsx
    warn   missing-deps  [hook:1]  var:userId  (line 5:2)  `userId` is used in this effect but not in its deps array, and its value may change between renders
       (1 trace step(s), rerun with --trace)
\end{sortie}

(Sorties de deux exécutions séparées, \code{check --no-color} sur chacun des
fichiers, juxtaposées ; les lignes de bilan sont omises.) On lira ces lignes en
détail à la \cref{sec:introduction-prise-en-main}. Remarquons seulement que les
deux diagnostics n'ont pas le même niveau : \code{error} pour le hook
conditionnel, \code{warn} pour la dépendance oubliée. Ce n'est pas un hasard,
et la \cref{sec:introduction-promesse} en donne la raison.

\subsection{Ce que la forme ne voit pas}

Reprenons maintenant la boucle infinie de \code{Counter}, mais écrite
\og proprement\fg{}, avec un tableau de deps complet :

\begin{lstlisting}[language=TypeScript,caption={\code{Counter} avec des deps complètes (\file{/tmp/ch01/deps_ok.tsx})},label={lst:introduction-deps-ok}]
import { useEffect, useState } from "react";

export function Counter() {
  const [n, setN] = useState(0);
  useEffect(() => {
    setN(n + 1);
  }, [n]);
  return <p>{n}</p>;
}
\end{lstlisting}

Pour \code{exhaustive-deps}, tout va bien : l'effet lit \code{n}, et \code{n}
est dans ses deps. Et pourtant le composant boucle exactement comme la
\cref{lst:introduction-loop} : à chaque exécution, l'effet écrit une valeur
différente de \code{n}, donc le dep \code{n} change, donc l'effet se
ré-exécute. La règle a vérifié une propriété de la \emph{forme} (\og la liste
est complète\fg{}) ; le défaut est une propriété de la \emph{suite des valeurs}
(\og \code{n} ne se stabilise jamais\fg{}). Aucun motif syntaxique ne peut
trancher ce second point en général : il faudrait savoir que \code{n + 1} n'est
jamais égal à \code{n}, et que rien ne vient arrêter la progression.
Comparez avec la variante bornée \code{if (n < 3) setN(n + 1)}, syntaxiquement
presque identique, qui s'arrête d'elle-même à \code{3}.

Même si l'on acceptait une règle approximative, du genre \og un effet qui
appelle le setter d'un état figurant dans ses deps\fg{}, il suffit d'une
indirection pour que le motif ne corresponde plus. Voici les trois indirections
que le plan de ce chapitre annonçait, chacune sur un programme de quelques
lignes.

\paragraph{Un alias de setter.} Le setter est rangé dans une autre variable
avant d'être appelé :

\begin{lstlisting}[language=TypeScript,caption={Alias de setter (\file{/tmp/ch01/alias.tsx})},label={lst:introduction-alias}]
import { useEffect, useState } from "react";

export function Counter() {
  const [n, setN] = useState(0);
  const bump = setN;
  useEffect(() => {
    bump(n + 1);
  }, [n]);
  return <p>{n}</p>;
}
\end{lstlisting}

Syntaxiquement, \code{bump} est une variable locale quelconque. Pour savoir
qu'appeler \code{bump} écrit l'état \code{n}, il faut suivre la \emph{valeur}
de \code{bump} depuis sa définition.

\paragraph{Un hook personnalisé.} L'état et l'effet vivent dans un hook
\code{useTicker}, que le composant se contente d'appeler :

\begin{lstlisting}[language=TypeScript,caption={Boucle cachée dans un hook personnalisé (\file{/tmp/ch01/custom_hook.tsx})},label={lst:introduction-custom-hook}]
import { useEffect, useState } from "react";

function useTicker() {
  const [n, setN] = useState(0);
  useEffect(() => {
    setN(n + 1);
  }, [n]);
  return n;
}

export function Clock() {
  const n = useTicker();
  return <p>{n}</p>;
}
\end{lstlisting}

Le défaut est dans \code{useTicker}, mais c'est \code{Clock} qui boucle : un
hook n'a pas de cycle de vie propre, il s'exécute \emph{dans} le composant qui
l'appelle. Si \code{useTicker} est dans un autre fichier, une règle qui
travaille fichier par fichier ne voit plus du tout l'effet depuis \code{Clock}.

\paragraph{Un setter passé en prop.} Le parent possède l'état et donne son
setter à un enfant, qui l'appelle dans un effet :

\begin{lstlisting}[language=TypeScript,caption={Setter du parent appelé par l'enfant (\file{/tmp/ch01/prop_setter.tsx})},label={lst:introduction-prop-setter}]
import { useEffect, useState } from "react";

function Child({ value, onChange }: { value: number; onChange: (v: number) => void }) {
  useEffect(() => {
    onChange(value + 1);
  }, [value, onChange]);
  return <p>{value}</p>;
}

export function Parent() {
  const [n, setN] = useState(0);
  return <Child value={n} onChange={setN} />;
}
\end{lstlisting}

Ici, aucun des deux composants n'est fautif isolément. \code{Parent} ne fait
que transmettre ; \code{Child} appelle une fonction \code{onChange} qu'il
reçoit, et ses deps sont complètes. La boucle n'existe que dans la
\emph{composition} : l'effet de \code{Child} écrit l'état de \code{Parent},
\code{Parent} se re-rend et re-rend \code{Child} avec une nouvelle
\code{value}, l'effet se ré-exécute.

\reactant{} signale les trois (une seule invocation sur les trois fichiers,
avec \code{--trace}) :

\begin{sortie}
  Child  (1 hooks)  prop_setter.tsx
    warn   cross-component-infinite-loop  [hook:0]  (line 4:2)  this effect calls `onChange`, a state setter of parent `Parent` (its deps do not provably gate it, so the effect can re-run every render). Parent re-renders → child re-renders → effect fires again: infinite loop
  Clock  (1 hooks)  custom_hook.tsx
    warn   infinite-loop  [hook:1]  (line 5:2)  this effect keeps pushing state `n` (its deps do not provably gate it, so the effect can re-run every render) to new values on every run. Potential infinite render loop
       → state `n` is written here [hook:2] (line 5:2)
       → the abstract value of state `n` kept growing and was widened at iteration 3
  Counter  (2 hooks)  alias.tsx
    warn   infinite-loop  [hook:0]  (line 6:2)  this effect keeps pushing state `n` (its deps do not provably gate it, so the effect can re-run every render) to new values on every run. Potential infinite render loop
       → state `n` is written here [hook:1] (line 6:2)
       → the abstract value of state `n` kept growing and was widened at iteration 3
   1 clean component(s) hidden, rerun with --show-clean

⚠  3 warning(s) across 3 file(s).
\end{sortie}

Trois détails de cette sortie annoncent la suite du manuscrit. Le finding de
\code{useTicker} est rapporté sur le composant \code{Clock}, qui est celui qui
boucle : le hook a été \emph{inliné} dans son appelant
(\cref{chap:splice-inlining}). Celui de \code{Child} nomme le parent
\code{Parent} dont il écrit l'état : l'analyse a suivi le setter à travers la
prop (\cref{chap:inter-composants}). Enfin, deux lignes parlent de la
\og valeur abstraite\fg{} de \code{n}, qui \og a grandi\fg{} puis a été
\og élargie à l'itération 3\fg{} : c'est le vocabulaire de l'interprétation
abstraite, que la \cref{sec:introduction-pari} présente.

Le hook personnalisé placé dans un autre fichier se comporte de même. Avec
\code{useTicker} déplacé dans \file{/tmp/ch01/xfile/useTicker.ts} et
\code{Clock} dans \file{/tmp/ch01/xfile/Clock.tsx}, l'analyse du dossier
entier donne :

\begin{sortie}
  Clock  (1 hooks)  Clock.tsx
    warn   infinite-loop  [hook:1]  (useTicker.ts:5:2)  this effect keeps pushing state `n` (its deps do not provably gate it, so the effect can re-run every render) to new values on every run. Potential infinite render loop
       → state `n` is written here [hook:2] (useTicker.ts:5:2)
       → the abstract value of state `n` kept growing and was widened at iteration 3

⚠  1 warning(s) across 2 file(s).
\end{sortie}

La position \code{(useTicker.ts:5:2)} désigne un autre fichier que celui du
composant : le finding est ancré dans le hook, mais porté par le composant qui
le subit (\adr{24}).

\begin{table}[htbp]\centering\small
\begin{tabularx}{\linewidth}{@{}lXX@{}}
  \toprule
  Programme & Ce que la forme montre & Ce qu'il faut calculer \\
  \midrule
  \code{Panel} (\cref{lst:introduction-cond}) & hook après un \code{return} conditionnel & rien de plus \\
  \code{Profile} (\cref{lst:introduction-deps}) & \code{userId} lu, absent des deps & que \code{userId} peut changer \\
  \code{Counter} (\cref{lst:introduction-deps-ok}) & deps complètes & la suite des valeurs de \code{n} \\
  alias (\cref{lst:introduction-alias}) & appel d'une variable locale & que \code{bump} \emph{est} le setter \\
  \code{Clock} (\cref{lst:introduction-custom-hook}) & appel d'un hook & le corps du hook, dans son appelant \\
  \code{Child}/\code{Parent} (\cref{lst:introduction-prop-setter}) & appel d'une prop & que la prop est le setter du parent \\
  \bottomrule
\end{tabularx}
\caption{Ce que la syntaxe montre et ce qu'il faut calculer pour conclure.}
\label{tab:introduction-syntaxe}
\end{table}

\Cref{tab:introduction-syntaxe} résume la leçon : dès qu'on quitte les deux
premières lignes, conclure demande de connaître des \emph{valeurs} (ce que
contient \code{n}, ce que désigne \code{bump} ou \code{onChange}) sur
\emph{toutes} les exécutions possibles.

\begin{remarque}
  Le \file{README.md} du dépôt situe \reactant{} par rapport à
  \code{eslint-plugin-react-hooks} et au React Compiler : le linter reste
  l'outil juste pour les règles lexicales des hooks, dans l'éditeur ; le
  compilateur mémoïse automatiquement et fait disparaître certains symptômes de
  références instables, sans corriger les défauts de logique. \reactant{} ne
  remplace ni l'un ni l'autre : il est plus lent (il calcule un point fixe sur
  tout le programme) et vise les défauts que la forme ne montre pas. Le README
  annonce par exemple 528 fichiers analysés en 1,8~s sur \code{excalidraw} et
  4\,195 fichiers en 76~s sur \code{dub} (\file{README.md}, l.~121--122 ;
  chiffres du README, non re-mesurés pour ce manuscrit).
\end{remarque}

\begin{piege}
  \og Le linter ne dit rien\fg{} signifie au mieux \og les deps sont
  complètes\fg{}, jamais \og l'effet se stabilise\fg{} : la
  \cref{lst:introduction-deps-ok} a un tableau de deps irréprochable et boucle.
  Inversement, compléter les deps pour faire taire le linter peut
  \emph{créer} la boucle. Partons de \code{useEffect(() => \{ setN(n + 1);
  \}, [])}, qui s'exécute une fois : \code{exhaustive-deps}, comme
  \regle{missing-deps}, réclame \code{n} (exercice~\ref{exo:introduction-once}) ;
  l'ajouter donne exactement la \cref{lst:introduction-deps-ok}.
\end{piege}

\section{Le pari : interpréter abstraitement le composant}
\label{sec:introduction-pari}

Si la forme ne suffit pas, que reste-t-il ? Exécuter le composant. Mais une
exécution ne prouve rien sur les autres : elle choisit une valeur de prop, un
ordre d'événements, un nombre de clics. Et on ne peut pas les essayer toutes :
une prop \code{number} a des milliards de valeurs, et une boucle de rendus peut
ne jamais s'arrêter. Le pari de \reactant{} est celui de l'\terme{interprétation
abstraite}~\cite{CousotCousot77} : exécuter le composant \emph{une fois}, mais
sur des \emph{descriptions d'ensembles de valeurs} plutôt que sur des valeurs.

\subsection{Exécuter sur des ensembles}

Revenons à \code{Counter} (\cref{lst:introduction-deps-ok}) et à sa variante
bornée :

\begin{lstlisting}[language=TypeScript,caption={\code{Bounded} : la même boucle, gardée (\file{/tmp/ch01/bounded.tsx})},label={lst:introduction-bounded}]
import { useEffect, useState } from "react";

export function Bounded() {
  const [n, setN] = useState(0);
  useEffect(() => {
    if (n < 3) setN(n + 1);
  }, [n]);
  return <p>{n}</p>;
}
\end{lstlisting}

Au lieu de dire \og \code{n} vaut \code{0}\fg{}, l'analyseur dit \og \code{n}
est un nombre entier compris entre \code{0} et \code{0}\fg{}, ce qu'on note par
l'intervalle $[0,0]$. Il rejoue alors un \emph{tour} de React : il interprète le
corps du composant (le rendu), puis le corps de l'effet. Dans \code{Counter},
l'effet écrit \code{n + 1}, soit un nombre de $[1,1]$. L'état \code{n} peut
donc valoir, à l'entrée du tour suivant, soit l'ancienne valeur, soit la
nouvelle : on réunit les deux en $[0,1]$. Au tour suivant l'effet écrit
$[1,2]$ et l'état devient $[0,2]$, puis $[0,3]$\dots{} La suite ne se stabilise
pas, et l'analyseur ne peut pas attendre. Il applique alors une opération
d'accélération, le \terme{widening} : la borne qui grandit est projetée d'un
coup vers une borne sûre. Pour \code{Counter}, cette borne est $+\infty$ ;
l'intervalle $[0,+\infty)$ est stable (l'effet y écrit $[1,+\infty)$, déjà
inclus), l'analyse s'arrête. Pour \code{Bounded}, la borne est le seuil
\code{3} lu dans la garde \code{n < 3} ; dans $[0,3]$, l'effet n'écrit que
lorsque $n<3$, donc des valeurs de $[1,3]$, et l'intervalle $[0,3]$ est lui
aussi stable.

\begin{figure}[htbp]\centering
\begin{tikzpicture}[node distance=5mm and 7mm,
    v/.style={bloc,minimum width=15mm,font=\small}]
  \node[etiquette,anchor=east] (lc) at (-0.4,0) {\code{Counter}};
  \node[v] (c0) at (0.6,0) {$[0,0]$};
  \node[v,right=of c0] (c1) {$[0,1]$};
  \node[v,right=of c1] (c2) {$[0,2]$};
  \node[v,right=17mm of c2,draw=ra-amber] (c3) {$[0,+\infty)$};
  \draw[fleche] (c0) -- node[etiquette,above]{tour} (c1);
  \draw[fleche] (c1) -- node[etiquette,above]{tour} (c2);
  \draw[flechemay] (c2) -- node[etiquette,above]{widening} (c3);
  \node[etiquette,right=3mm of c3,align=left] {stable, non borné\\$\Rightarrow$ \regle{infinite-loop}};
  \node[etiquette,anchor=east] (lb) at (-0.4,-1.3) {\code{Bounded}};
  \node[v] (b0) at (0.6,-1.3) {$[0,0]$};
  \node[v,right=of b0] (b1) {$[0,1]$};
  \node[v,right=of b1] (b2) {$[0,2]$};
  \node[v,right=17mm of b2] (b3) {$[0,3]$};
  \draw[fleche] (b0) -- node[etiquette,above]{tour} (b1);
  \draw[fleche] (b1) -- node[etiquette,above]{tour} (b2);
  \draw[flechemay] (b2) -- node[etiquette,above]{widening} (b3);
  \node[etiquette,right=3mm of b3,align=left] {stable, borné par\\la garde $\Rightarrow$ silence};
\end{tikzpicture}
\caption{Valeur abstraite de l'état \code{n} au fil des tours abstraits (déroulé
simplifié ; le détail des itérations et le seuil du widening sont au
\cref{chap:point-fixe-composant}). Les valeurs finales $[0,+\infty)$ et
$[0,3]$ sont celles que le moteur calcule effectivement au commit
\snapshotCommit{}.}
\label{fig:introduction-intervalles}
\end{figure}
% Valeurs finales vérifiées par une sonde hors dépôt (/tmp/ch01probe, appelle
% analyze_program sur /tmp/ch01/{deps_ok,bounded}.tsx) : Counter state[0] =
% number[0, inf], Bounded state[0] = number[0, 3], 3 itérations, slot 0 élargi
% dans les deux cas. Les états intermédiaires [0,1], [0,2] sont un déroulé à la
% main (cf. notes/16 Ex. 4), non tracés pas à pas.

\Cref{fig:introduction-intervalles} résume ces deux déroulés. Ils ont un point
commun essentiel : le résultat final \emph{contient} toutes les valeurs que
\code{n} peut réellement prendre. Dans \code{Counter}, \code{n} prend les
valeurs $0,1,2,\dots$, toutes dans $[0,+\infty)$ ; dans \code{Bounded}, les
valeurs $0,1,2,3$, toutes dans $[0,3]$. C'est cette inclusion qui rend le
raisonnement valable : l'analyse \emph{sur-approxime}, elle peut se tromper en
imaginant des valeurs impossibles, jamais en oubliant des valeurs possibles.

\begin{definition}{Valeur abstraite, point fixe (informel)}{introduction-abstraction}
  Une \terme{valeur abstraite} décrit un ensemble de valeurs concrètes (un
  intervalle décrit les nombres qu'il contient). L'analyse d'un composant
  répète des tours abstraits (rendu, effets, handlers), en réunissant à chaque
  tour ce que l'état pouvait valoir et ce qui vient d'y être écrit, jusqu'à ce
  qu'un tour n'ajoute plus rien : c'est le \terme{point fixe}. Le widening
  garantit qu'on l'atteint en un nombre fini de tours. Les définitions exactes
  (treillis, join, concrétisation, widening) sont au
  \cref{chap:interpretation-abstraite}.
\end{definition}

Le point fixe atteint, les règles n'ont plus qu'à \emph{lire} le résultat.
Pour \code{Counter}, l'état a dû être élargi et l'effet y écrit sans borne :
c'est la signature d'une boucle de rendus, et \regle{infinite-loop} la
signale. Pour \code{Bounded}, l'état a aussi été élargi (le widening est une
approximation, l'analyseur le dit sous \code{--info}), mais l'écriture reste
bornée par la garde : pas de boucle, et la règle l'atteste.

\begin{sortie}
  Bounded  (2 hooks)  fixed.tsx
    info   widening-info  (line 26:2)  state `n` kept changing during analysis and was approximated to converge, so findings that depend on it may be imprecise
       (2 trace step(s), rerun with --trace)
    verified  always-unstable-deps  no deps array is defeated by an always-fresh reference
    verified  conditional-hook  all hooks run unconditionally, in a stable order
    verified  derived-state  no effect merely mirrors other state
    verified  infinite-loop  no effect diverges into an infinite render loop
    ...
\end{sortie}

(Extrait de \code{check --no-color --info --show-clean} sur
\file{/tmp/ch01/fixed.tsx}, qui regroupe les versions corrigées des exemples
du chapitre ; les six autres lignes \code{verified} sont élaguées.)

\subsection{Pas seulement des nombres}

Les intervalles ne suffisent pas pour React. Les trois défauts de la
\cref{sec:introduction-trois-bugs} demandent chacun une information différente
sur les valeurs, et \reactant{} les abstrait toutes à la fois.

\begin{itemize}
  \item \textbf{La valeur.} \code{n + 1} diffère-t-il de \code{n} ? Le domaine
    des valeurs combine un intervalle pour les nombres, des ensembles bornés de
    chaînes, des booléens, \code{null}, \code{undefined}, dans un produit
    (\cref{chap:treillis-simples,chap:statevalue-stabilite}).
  \item \textbf{L'identité.} React compare les deps et les états avec
    \code{Object.is}, qui teste l'identité des objets et non leur contenu. Un
    objet littéral \code{\{\ldots\}} est une référence \emph{neuve} à chaque
    évaluation. Le domaine suit donc aussi la \emph{stabilité} d'une valeur :
    stable, changeant seulement quand un état change, ou neuve à chaque rendu
    (\cref{chap:statevalue-stabilite}).
  \item \textbf{Le moment.} Qui écrit quel état, pendant le rendu, dans un
    effet, ou plus tard dans un callback (\code{setInterval}, \code{.then}) ?
    Avec quelle valeur capturée ? Ces faits sont calculés une fois par le moteur
    sous forme de \emph{relations} (\cref{chap:relations}).
\end{itemize}

\begin{exemple}{L'identité, pas la valeur}{introduction-churn}
  L'effet suivant ne \og change\fg{} rien au contenu de \code{query} dès le
  deuxième passage (\code{trim} d'une chaîne déjà nettoyée), et pourtant il
  boucle :
\begin{lstlisting}[language=TypeScript,caption={Un objet recréé à chaque exécution de l'effet (\file{/tmp/ch01/churn.tsx})},label={lst:introduction-churn}]
import { useEffect, useState } from "react";

export function Filters() {
  const [query, setQuery] = useState({ term: "" });
  useEffect(() => {
    setQuery({ ...query, term: query.term.trim() });
  }, [query]);
  return <input value={query.term} />;
}
\end{lstlisting}
  Le littéral \code{\{ ...query, term: ... \}} est un objet neuf : pour
  \code{Object.is}, le dep \code{query} a changé, l'effet repart. L'analyse le
  sait sans rien connaître du contenu des chaînes :
\begin{sortie}
  Filters  (2 hooks)  /tmp/ch01/churn.tsx
    error  infinite-loop  [hook:0]  (line 5:2)  this effect recreates object state `query` it depends on. Every run stores a fresh reference (`Object.is` always fails) and re-triggers itself: infinite render loop
       (1 trace step(s), rerun with --trace)
\end{sortie}
  Notons le niveau : \code{error}, alors que la boucle de \code{Counter}
  n'était qu'un \code{warn}. La \cref{sec:introduction-promesse} explique cette
  différence.
\end{exemple}

\subsection{Quel React on abstrait}

Pour que \og sur-approximer\fg{} ait un sens, il faut savoir ce qu'on
sur-approxime : une \emph{sémantique concrète} de React, qui dit précisément
ce qu'une exécution peut faire. React n'en publie pas ; sa documentation décrit
le comportement en prose~\cite{ReactDocs}. \reactant{} a fait le choix
d'adosser son modèle à une sémantique formelle publiée.

\begin{decision}[ADR-001]
  \textbf{Problème.} Sans sémantique concrète explicite, on ne peut ni établir
  la soundness de l'analyseur ni écrire ses fonctions de transfert autrement
  qu'à l'intuition. \textbf{Décision.} Adopter comme sémantique concrète
  React-tRace (Lee, Ahn et Yi, OOPSLA 2025), dont les étapes \emph{init},
  \emph{effect} et \emph{check} correspondent à une itération du point fixe.
  \textbf{Limite acceptée.} React-tRace ne couvre qu'un noyau
  (\code{useState}, \code{useEffect}) ; les extensions du projet (deps,
  mémos, contexte, props\dots) n'ont pas de garantie formelle équivalente.
  Le \cref{chap:react-modele} présente cette sémantique et ce que le projet
  en garde (\adr{1}).
\end{decision}
% BIB: React-tRace — clé LeeAhnYi2025 proposée au chapitre 02
% (@article{LeeAhnYi2025, ..., doi={10.1145/3763067}, note={arXiv:2507.05234}}).

Concrètement, le tour abstrait que rejoue le moteur est celui de la
\cref{fig:introduction-cycle-react}, légèrement grossi pour rester sûr : à
chaque tour, il interprète le rendu, les mémos, \emph{tous} les effets (qu'ils
soient ou non re-déclenchés, ce qui couvre toutes les combinaisons de deps) et
les handlers (qu'un utilisateur pourrait déclencher à tout moment), puis réunit
les états écrits avec les états précédents. Le \cref{chap:point-fixe-composant}
décrit cette boucle, et en quoi chaque grossissement préserve la
sur-approximation.

\subsection{Le pipeline}

Interpréter un composant demande d'abord de le lire. \reactant{} n'interprète
pas directement le texte source : il le traduit en une représentation
intermédiaire (IR) conçue pour l'analyse, un graphe de flot de contrôle par
composant, où les hooks sont des objets de première classe. La
\cref{fig:introduction-pipeline} montre les étapes d'une invocation, avec le
module Rust qui porte chacune et le type qui passe à l'étape suivante.

\begin{figure}[htbp]\centering
\begin{tikzpicture}[node distance=6mm and 13mm,
    etape/.style={bloc,text width=20mm,minimum height=13mm},
    type/.style={etiquette,font=\scriptsize\ttfamily,text=ra-gray}]
  \node[bloc,text width=13mm] (src) {fichiers\\\code{.tsx}};
  \node[etape,right=9mm of src] (parse) {parse (oxc)\\{\scriptsize\file{src/resolver}}};
  \node[etape,right=of parse] (lower) {lowering\\{\scriptsize\file{src/lowering}}};
  \node[etape,right=of lower] (ir) {IR\\{\scriptsize\file{src/ir}}};
  \node[etape,below=18mm of ir] (engine) {moteur\\{\scriptsize\file{src/engine}}\\{\scriptsize\file{src/domains}}};
  \node[etape,left=of engine] (rel) {relations\\{\scriptsize\file{src/engine}}};
  \node[etape,left=of rel] (rules) {règles\\{\scriptsize\file{src/rules}}};
  \node[etape,left=of rules] (driver) {driver\\{\scriptsize\file{src/driver}}\\{\scriptsize\file{src/cli}}};
  \node[bloc,below=12mm of driver,text width=22mm] (out) {sortie humaine\\ou JSON};
  \draw[fleche] (src) -- (parse);
  \draw[fleche] (parse) -- node[type,above]{AST} (lower);
  \draw[fleche] (lower) -- node[type,above]{} (ir);
  \draw[fleche] (ir) -- node[type,right]{ComponentIR} (engine);
  \draw[fleche] (engine) -- node[type,below,align=center,yshift=-1mm]{Analysis-\\Result} (rel);
  \draw[fleche] (rel) -- node[type,below,align=center,yshift=-1mm]{Program-\\Relations} (rules);
  \draw[fleche] (rules) -- node[type,below,yshift=-1mm]{Diagnostic} (driver);
  \draw[fleche] (driver) -- node[type,right]{CheckReport} (out);
\end{tikzpicture}
\caption{Le pipeline d'une invocation \code{reactant check}. Chaque boîte est
une couche du code ; sur les flèches, le type principal qui passe d'une couche
à la suivante. Le \cref{chap:tour-architecture} détaille chaque étape et ses
frontières.}
\label{fig:introduction-pipeline}
\end{figure}

On y lit trois décisions qui structurent tout le projet.
\begin{enumerate}
  \item \textbf{Après le lowering, plus aucun AST.} Le front-end utilise le
    parseur \code{oxc}~\cite{oxc}, détecte les composants, hooks et
    utilitaires, et produit un \code{ComponentIR} par composant ; le moteur ne
    voit jamais la syntaxe d'origine (\adr{3},
    \cref{chap:frontend-detection,chap:ir,chap:lowering-cfg}).
  \item \textbf{Le moteur calcule, les règles lisent.} Le moteur exécute le
    point fixe de chaque composant et dérive des relations (qui écrit quel
    état, quand, avec quelle valeur ; quels effets se re-déclenchent) ; les
    règles ne parcourent pas l'IR, elles interrogent ces résultats (\adr{42},
    \cref{chap:relations,chap:api-regles}).
  \item \textbf{L'analyse est de programme.} Un parent est analysé avec ses
    enfants, les hooks personnalisés et utilitaires sont inlinés dans leurs
    appelants, à travers les fichiers (\cref{chap:splice-inlining,chap:inter-composants}).
    C'est ce qui a permis de voir les trois indirections de la
    \cref{sec:introduction-linter}.
\end{enumerate}

\section{La promesse : trois niveaux de diagnostic et la soundness}
\label{sec:introduction-promesse}

Rassemblons les niveaux observés jusqu'ici :
\begin{center}\small
\begin{tabular}{@{}lll@{}}
  \toprule
  Programme & Diagnostic & Niveau \\
  \midrule
  \code{Panel} (\cref{lst:introduction-cond}) & \regle{conditional-hook} & \Error \\
  \code{Filters} (\cref{lst:introduction-churn}) & \regle{infinite-loop} & \Error \\
  \code{Timer} (\cref{lst:introduction-timer}) & \regle{stale-closure} & \Error \\
  \code{Timer} (\cref{lst:introduction-timer}) & \regle{missing-deps} & \Warning \\
  \code{Profile} (\cref{lst:introduction-deps}) & \regle{missing-deps} & \Warning \\
  \code{Counter} (\cref{lst:introduction-loop,lst:introduction-deps-ok}) & \regle{infinite-loop} & \Warning \\
  \code{Bounded} (\cref{lst:introduction-bounded}) & \regle{widening-info} & \Info \\
  \bottomrule
\end{tabular}
\end{center}
(La sortie de \code{Timer} est commentée à la \cref{sec:introduction-prise-en-main}.)
Le niveau n'est pas une impression de gravité : c'est une promesse sur la
\emph{nature} de ce que l'outil sait.

\subsection{Trois niveaux, trois promesses}

\begin{definition}{Niveaux de diagnostic}{introduction-niveaux}
  \begin{itemize}
    \item \Error{} : le défaut est \emph{certain} dès que le code signalé
      s'exécute. L'outil détient une preuve de \emph{toute} la conclusion,
      pas seulement d'une de ses conditions.
    \item \Warning{} : le défaut est \emph{possible} (il dépend d'un chemin,
      d'une entrée, ou l'analyse a sur-approximé), ou bien un fait est certain
      mais son coût ne l'est pas (une référence neuve à chaque rendu, un rendu
      gaspillé).
    \item \Info{} : une limite de l'analyse (un widening, un hook qu'elle n'a
      pas pu lire) ou un motif qui a l'air voulu. Masqué par défaut, affiché
      avec \code{--info}.
  \end{itemize}
\end{definition}

Cette définition n'est pas seulement une convention de documentation : elle
est écrite dans le type qui porte le niveau, et c'est la même que celle des
invariants du projet dans \file{CLAUDE.md}.

\rustsnippet[label={lst:introduction-severity}]{src/rules/api/diagnostic.rs}{23}{38}{les trois niveaux et leur contrat}

Relisons nos exemples avec cette grille.
\begin{itemize}
  \item \code{Profile} : le défaut n'existe que si \code{userId} change
    pendant la vie du composant ; rien, dans le programme analysé, ne dit
    qu'il change ni qu'il ne change pas. Défaut possible : \Warning.
  \item \code{Panel} : le \code{useState} n'est pas sur tous les chemins du
    rendu (le retour anticipé l'évite). C'est un fait de structure, prouvable
    en regardant le graphe de flot : aucun chemin de l'entrée à la sortie ne
    peut être oublié. \Error.
  \item \code{Filters} : trois faits certains se composent. L'effet écrit
    \code{query} sur tous ses chemins ; il y écrit toujours un objet neuf ;
    \code{query} est exactement l'un de ses deps. La conclusion \og il se
    re-déclenche lui-même à chaque exécution\fg{} est entièrement prouvée.
    \Error.
  \item \code{Counter} : la boucle est pourtant tout aussi certaine. Mais ce
    que l'analyse sait, c'est que l'intervalle de \code{n} a dû être élargi et
    que l'effet y écrit sans borne. Un widening peut aussi venir d'une simple
    imprécision sur une boucle qui converge ; il ne constitue pas une preuve.
    L'outil ne dispose pas, au commit \snapshotCommit{}, d'une primitive qui
    prouverait \og la valeur écrite diffère toujours de la valeur
    courante\fg{} : il s'en tient à \Warning{} (\issue{144}).
\end{itemize}

Le cas \code{Counter} enseigne l'asymétrie essentielle : un niveau trop
\emph{bas} est permis, un niveau trop \emph{haut} ne l'est pas. Un \Warning{}
sur un défaut certain est une perte de précision ; une \Error{} sur un défaut
seulement possible est une promesse rompue.

\subsection{Une Error se construit à partir d'une preuve}

Comment le code empêche-t-il une règle de promettre trop ? Par le typage. Le
constructeur d'une \Error{} exige un argument de type \code{Certified<E>},
un \og jeton de preuve\fg{} que seules les primitives de preuve du moteur
savent fabriquer (\cref{chap:api-regles}) ; un \Warning{} ou une \Info{} se
construisent librement.

\rustsnippet[label={lst:introduction-error-ctor}]{src/rules/api/diagnostic.rs}{157}{191}{le constructeur d'une \Error{}, contre les constructeurs libres d'un \Warning{} et d'une \Info{}}

Une règle qui n'a qu'un fait may (vrai sur certains chemins seulement) n'a
pas de \code{Certified} à passer : écrire une \Error{} à partir de lui est une
erreur de compilation, pas une faute de relecture.

\begin{decision}[ADR-021]
  \textbf{Problème.} Un faux négatif réel s'était glissé dans un utilitaire
  partagé des règles, et n'importe quel booléen pouvait devenir une
  \code{Severity::Error}. \textbf{Décision.} Faire de la \emph{polarité}
  d'un fait (vrai sur tous les chemins, sur certains, ou inconnu) un type :
  \code{MustResult}, \code{May}, et un jeton \code{Certified} que seules les
  primitives \code{must_*} produisent ; \code{Diagnostic::error} n'accepte que
  ce jeton. Les règles existantes ont été migrées d'un bloc plutôt que
  progressivement, pour ne pas laisser de fenêtre ouverte aux faux négatifs.
  \textbf{Précision ultérieure (\issue{142}).} La preuve doit couvrir
  \emph{toute} la conclusion : une Error de \regle{stale-closure} avait été
  produite à partir d'une preuve ne couvrant qu'une de ses trois conditions,
  ce qui a conduit à resserrer la règle et la doctrine des niveaux
  (\adr{21}, \cref{chap:api-regles,chap:histoire-doctrine}).
\end{decision}

\subsection{La soundness : pas de faux négatif}

Reste la promesse la plus forte, qui porte sur le \emph{silence}. Parlons
d'abord vocabulaire.

\begin{definition}{Faux positif, faux négatif, soundness (informel)}{introduction-soundness}
  Un \terme{faux positif} (FP) est un diagnostic sur un défaut qui n'existe
  pas ; un \terme{faux négatif} (FN) est un défaut réel que l'outil ne signale
  pas. Une analyse est \terme{sound} (on parle de sa \emph{soundness}) si elle
  calcule un sur-ensemble des comportements possibles du programme : tout ce
  qui peut arriver est dans ce qu'elle a calculé. Conséquence : ses FN sont
  exclus, ses FP tolérés. La définition exacte est au
  \cref{chap:interpretation-abstraite}.
\end{definition}

La soundness est l'invariant numéro un du projet (\file{CLAUDE.md}) : \og faux
positifs tolérés, faux négatifs INTERDITS\fg{}. Elle se lit à plusieurs
endroits de la sortie.
\begin{itemize}
  \item \textbf{Le silence est une affirmation.} Quand une règle applicable
    n'a rien trouvé, \code{--info} l'écrit : \code{verified infinite-loop
    no effect diverges into an infinite render loop}. Cette ligne n'est
    publiée que si l'analyse du composant n'a pas été tronquée.
  \item \textbf{Ce qui n'a pas été lu est dit.} Si l'on analyse
    \file{/tmp/ch01/xfile/Clock.tsx} seul, sans le fichier du hook qu'il
    importe, l'outil ne conclut pas \og aucun problème\fg{} :
\begin{sortie}
   1 clean component(s) hidden, rerun with --show-clean

⚠  1 file(s), no findings, but parts of this run were not analyzed, so this is not a clean bill.
   not analyzed:
     • 1 imported file(s) resolved outside the analysed set and were never read. Pass them on the command line to analyse them (useTicker.ts)
\end{sortie}
    Le code de sortie reste \code{0} (aucun diagnostic), mais la phrase finale
    refuse le \og ✓ no issues found\fg{} : on ne délivre pas de certificat de
    propreté pour du code qu'on n'a pas lu (\cref{chap:projet-cli-driver}).
  \item \textbf{L'inconnu est ⊤, pas \og rien\fg{}.} Un hook de bibliothèque
    que l'outil ne sait pas modéliser renvoie une valeur \emph{inconnue} (le
    top du treillis, ⊤), qui peut être n'importe quoi. Les règles qui la lisent
    doivent alors supposer le pire, et une \Info{} \regle{analysis-limit}
    indique la limite (\cref{chap:regles-effets}).
\end{itemize}

\begin{piege}
  La soundness est un \emph{invariant visé}, pas un théorème démontré sur le
  code. La page \file{docs/limitations.md} du dépôt recense des
  limites connues, et le tracker du projet porte des issues étiquetées
  \code{soundness-bug} ouvertes au commit \snapshotCommit{} (par exemple
  \issue{157}, \issue{143}, \issue{7}). La rédaction de ce manuscrit a aussi
  mis au jour des faux négatifs. Le plus simple tient en six lignes :
\begin{lstlisting}[language=TypeScript,caption={Une boucle certaine que la CLI ne signale pas (\file{/tmp/ch01/updater_dep.tsx})},label={lst:introduction-fn-updater}]
import { useEffect, useState } from "react";

export function D() {
  const [count, setCount] = useState(0);
  useEffect(() => {
    setCount((c) => c + 1);
  }, [count]);
  return <p>{count}</p>;
}
\end{lstlisting}
  L'updater fait croître \code{count} à chaque exécution, \code{count} est un
  dep, l'effet boucle. Pourtant :
\begin{sortie}
   1 clean component(s) hidden, rerun with --show-clean

✓  1 file(s) no issues found.
\end{sortie}
  La cause (le chemin \og programme\fg{} du moteur écrit la fonction
  \code{(c) => c + 1} elle-même dans l'état, au lieu d'appliquer l'updater)
  est analysée au \cref{chap:limites-perspectives}. Le lecteur doit donc
  retenir deux choses à la fois : la conception entière vise l'absence de FN,
  et les écarts à cette visée sont des défauts à corriger à la racine, pas des
  comportements acceptés.
\end{piege}
% Défaut documenté dans notes/16 §8.1 (sonde : PROGRAM state[0] = ⊤, aucun
% widening ; les tests intra de tests/functional_updater.rs passent). Pas
% d'issue ouverte au moment de la rédaction (gh issue list, 2026-09-28).

\section{Première prise en main}
\label{sec:introduction-prise-en-main}

Il est temps de lancer l'outil soi-même. Toutes les sorties de ce manuscrit
ont été produites au commit \snapshotCommit{}, depuis la racine du dépôt, par
le binaire de développement.

\subsection{Construire et lancer}

\begin{shell}
cargo build --release          # binaire : target/release/reactant
cargo run -q -- check /tmp/ch01/timer.tsx
cargo run -q -- check src/     # un dossier : parcours récursif des .ts/.tsx/.js/.jsx
\end{shell}

\code{check} est la sous-commande par défaut (\code{reactant src/} équivaut à
\code{reactant check src/}). Sans chemin, c'est le dossier courant qui est
analysé. Les autres sous-commandes sont \code{rules} (le catalogue),
\code{explain <rule>} (la documentation d'un diagnostic, avec un exemple et
une correction), \code{schemas} (les schémas JSON de la configuration et des
packs). Côté JavaScript, le même analyseur est distribué sur npm, compilé en
WebAssembly : \code{npx reactant-analyzer} (\cref{chap:distribution}).

\subsection{Lire la sortie humaine}

Lançons l'analyse sur \code{Timer} (\cref{lst:introduction-timer}), en
demandant les chaînes causales avec \code{--trace}. On se place dans
\file{/tmp/ch01} pour que les chemins affichés soient courts.

\begin{sortie}
  Timer  (2 hooks)  timer.tsx
    warn   missing-deps  [hook:1]  var:seconds  (line 5:2)  `seconds` is used in this effect but not in its deps array, and it is recreated on every render
       → `seconds` is read here [hook:1] (line 5:2)
    error  stale-closure  [hook:1]  var:seconds  (line 6:10)  the `setInterval` callback registered by this mount-only effect reads `seconds` and writes it back. `seconds` was captured once at mount, so every firing recomputes from the same frozen value and the state can never advance past its first update
       → `setInterval` has side effects (subscriptions/requests/timers re-fire on every call) [hook:1] (line 6:10)
       → `seconds` is captured at registration time, so the callback keeps this value, not the latest one [hook:0] (line 6:10)
       → state `seconds` is written here [hook:0] (line 7:6)

⚠  1 error(s), 1 warning(s) across 1 file(s).
\end{sortie}

Décomposons.
\begin{itemize}
  \item \textbf{L'en-tête} \code{Timer (2 hooks) timer.tsx} nomme le composant,
    le nombre de hooks que l'IR lui connaît et son fichier. Les hooks sont
    numérotés dans l'ordre d'appel, à partir de 0 : \code{[hook:0]} est le
    \code{useState}, \code{[hook:1]} le \code{useEffect}. Ce numéro est le
    \emph{label} du hook, celui qui identifie un slot d'état pour React comme
    pour l'analyseur (\cref{chap:ir}).
  \item \textbf{Chaque diagnostic} donne son niveau (\code{error},
    \code{warn}, \code{info}), le nom de la règle, le hook concerné, la
    variable en cause (\code{var:seconds}), la position \code{(line L:C)}
    --- ligne comptée à partir de 1, colonne à partir de 0 --- et un message.
    Ici, \code{(line 6:10)} pointe sur l'appel \code{setInterval(...)}.
  \item \textbf{Les lignes \og →\fg{}} sont la \emph{chaîne de témoins} du
    diagnostic : les faits, typés et localisés, qui ont conduit à la
    conclusion (\adr{19}, \cref{chap:api-regles}). Pour
    \regle{stale-closure}, ce sont exactement les trois ingrédients du défaut
    raconté à la \cref{sec:introduction-trois-bugs} : \code{setInterval}
    rappellera le callback ; le callback a capturé \code{seconds} au moment
    de l'enregistrement ; il réécrit l'état \code{seconds}. Sans
    \code{--trace}, ces lignes sont remplacées par un rappel
    \code{(3 trace step(s), rerun with --trace)}.
  \item \textbf{La dernière ligne} fait le bilan. Le code de sortie est
    \code{1} : par défaut (\code{--fail-on warning}), une \Error{} ou un
    \Warning{} fait échouer la commande, ce qui la rend utilisable telle quelle
    en intégration continue.
\end{itemize}

Les deux diagnostics racontent le même défaut à deux niveaux de preuve.
\regle{missing-deps} constate une propriété locale (\code{seconds} est lu
et n'est pas déclaré), en parité avec \code{exhaustive-deps} : défaut
possible, \Warning. \regle{stale-closure} prouve la \emph{conséquence} : avec
des deps \code{[]} et un callback répétitif qui réécrit ce qu'il lit, l'état
ne peut jamais dépasser sa première mise à jour. C'est ce que l'on a constaté
à la main (l'affichage bloqué à \og 1 s\fg{}), et la preuve couvre toute la
conclusion : \Error{} (\cref{chap:regles-effets}).

\begin{piege}
  La formulation du motif ajouté par \regle{missing-deps} (\og and it is
  recreated on every render\fg{}, \og and its value may change between
  renders\fg{}) n'est pas un jugement de l'auteur de la règle : elle traduit
  en anglais courant la valeur abstraite que le moteur a calculée pour la
  variable (\code{describe_value}, \srcline{src/rules/helpers/mod.rs}{42}).
  Elle peut donc surprendre, comme ici pour un nombre : \code{seconds} n'a
  pas la forme d'une référence, mais son intervalle ne s'est jamais réduit à
  un point pendant l'analyse (\cref{sec:introduction-pari}), ce qui suffit à
  le classer \og recreated on every render\fg{} au même titre qu'un objet
  neuf. Les \cref{chap:statevalue-stabilite,chap:regles-effets}
  expliquent cette traduction.
\end{piege}

La correction est celle de la \cref{sec:introduction-trois-bugs} : un updater,
\code{setSeconds((s) => s + 1)}. Le fichier \file{/tmp/ch01/fixed.tsx}
rassemble les versions corrigées de \code{Timer}, de \code{Profile} (avec
\code{[userId]}) et de \code{Bounded}. Sans option, la sortie est :

\begin{sortie}
   3 clean component(s) hidden, rerun with --show-clean

✓  1 file(s) no issues found.
\end{sortie}

Avec \code{--info --show-clean}, chaque composant affiche aussi les
vérifications applicables qui sont passées ; pour \code{Timer}, on lit entre
autres :

\begin{sortie}
  Timer  (2 hooks)  fixed.tsx  ✓
    verified  conditional-hook  all hooks run unconditionally, in a stable order
    ...
    verified  stale-closure  no long-lived callback captures a stale state value
    ...
\end{sortie}

Le signe ✓ après le fichier marque un composant sans diagnostic.

\begin{table}[htbp]\centering\small
\begin{tabularx}{\linewidth}{@{}lX@{}}
  \toprule
  Option & Effet \\
  \midrule
  \code{--trace} & affiche la chaîne de témoins de chaque diagnostic \\
  \code{--info} & affiche les \Info{} et les lignes \code{verified} / \code{suspended} \\
  \code{--show-clean} & affiche aussi les composants sans diagnostic \\
  \code{--format json} & un unique document JSON sur la sortie standard \\
  \code{--fail-on error|warning|never} & seuil du code de sortie \code{1} (défaut \code{warning}) \\
  \code{--rule}, \code{--ignore-rule} & ne garder / retirer qu'un nom de diagnostic \\
  \code{--follow-imports} & lire aussi les fichiers importés, transitivement \\
  \code{--verbose} & statistiques du moteur sur la sortie d'erreur \\
  \code{--no-color} & sortie sans couleurs ANSI \\
  \bottomrule
\end{tabularx}
\caption{Les options de \code{check} utilisées dans ce manuscrit (extrait de
\file{docs/usage.md} ; la liste complète et la configuration sont au
\cref{chap:projet-cli-driver}).}
\label{tab:introduction-options}
\end{table}

\Cref{tab:introduction-options} rassemble les options les plus utiles.
\code{--verbose} est une fenêtre sur le moteur : sur la boucle de la
\cref{lst:introduction-loop}, elle affiche, en plus du diagnostic,
\begin{sortie}
[verbose] symbol graph: 1 nodes, topo order = [Counter@loop.tsx]
[verbose] 1 components analyzed
[verbose] cache hits: 0, misses: 0
  [verbose] Counter: 3 iteration(s), widened: [0]
\end{sortie}
c'est-à-dire trois tours abstraits, et le slot de label 0 (\code{n}) élargi,
comme sur la \cref{fig:introduction-intervalles}.

\subsection{Lire la sortie JSON}

Pour une machine (intégration continue, éditeur, agent), \code{--format json}
écrit un document unique au schéma versionné. Voici celui de \code{Timer},
où le premier diagnostic (\regle{missing-deps}) a été élagué :

\begin{lstlisting}[language=JSON,caption={Sortie \code{check --format json timer.tsx} (élaguée : le diagnostic \regle{missing-deps} est remplacé par \code{...})},label={lst:introduction-json}]
{
  "version": 2,
  "files_analyzed": 1,
  "parse_errors": [],
  "diagnostics": [
    ...,
    {
      "rule": "stale-closure",
      "severity": "error",
      "component": "Timer",
      "file": "timer.tsx",
      "component_file": "timer.tsx",
      "line": 6,
      "col": 10,
      "hook_label": 1,
      "var": "seconds",
      "message": "the `setInterval` callback registered by this mount-only effect reads `seconds` and writes it back. `seconds` was captured once at mount, so every firing recomputes from the same frozen value and the state can never advance past its first update",
      "notes": [
        {
          "message": "`setInterval` has side effects (subscriptions/requests/timers re-fire on every call)",
          "kind": "call",
          "hook_label": 1,
          "file": "timer.tsx",
          "line": 6,
          "col": 10,
          "callee": "setInterval",
          "effect_class": "effectful"
        },
        {
          "message": "`seconds` is captured at registration time, so the callback keeps this value, not the latest one",
          "kind": "capture",
          "hook_label": 0,
          "file": "timer.tsx",
          "line": 6,
          "col": 10,
          "what": "seconds"
        },
        {
          "message": "state `seconds` is written here",
          "kind": "write",
          "hook_label": 0,
          "file": "timer.tsx",
          "line": 7,
          "col": 6,
          "slot": 0,
          "value_class": "unknown"
        }
      ]
    }
  ],
  "blind_spots": [],
  "summary": {
    "errors": 1,
    "warnings": 1,
    "infos": 0,
    "components_analyzed": 1,
    "exit_code": 1
  }
}
\end{lstlisting}

On retrouve champ pour champ la sortie humaine, plus quelques informations
structurées.
\begin{itemize}
  \item \code{version} est la version du schéma (2) ; \code{files_analyzed}
    et \code{parse_errors} disent ce qui a été lu, et si le parseur a dû
    récupérer d'erreurs de syntaxe.
  \item Dans chaque diagnostic, \code{file} est le fichier où pointe la
    position, \code{component_file} celui où le composant est défini ; les
    deux diffèrent quand le finding est ancré dans un hook inliné depuis un
    autre fichier (\cref{sec:introduction-linter}, exemple \code{Clock}).
    \code{line} commence à 1, \code{col} à 0.
  \item \code{notes} est la chaîne de témoins ; chaque étape porte un
    \code{kind} pris dans un vocabulaire fermé (\code{call}, \code{capture},
    \code{write}, \code{read}, \code{widen}\dots) et les champs propres à ce
    genre (\code{callee}, \code{what}, \code{slot}\dots).
  \item \code{blind_spots} liste ce que l'exécution sait ne pas avoir lu. Non
    vide, il signifie que les compteurs du \code{summary} sont une borne
    inférieure, pas un verdict ; c'est la version machine de la phrase
    \og this is not a clean bill\fg{} vue plus haut.
  \item \code{summary.exit_code} reproduit le code de sortie du processus.
\end{itemize}

\begin{table}[htbp]\centering\small
\begin{tabular}{@{}cl@{}}
  \toprule
  Code & Signification \\
  \midrule
  \code{0} & aucun diagnostic au niveau du seuil \code{--fail-on} ou au-dessus \\
  \code{1} & au moins un diagnostic au niveau du seuil ou au-dessus \\
  \code{2} & erreur d'usage ou d'entrée : aucun fichier, chemin inexistant, règle inconnue \\
  \bottomrule
\end{tabular}
\caption{Codes de sortie de \code{reactant check} (\file{docs/usage.md}). Les
\Info{} et les \code{blind_spots} ne changent jamais le code de sortie.}
\label{tab:introduction-exit}
\end{table}

\Cref{tab:introduction-exit} donne les codes de sortie. Enfin,
\code{reactant explain <rule>} affiche la documentation d'un diagnostic. Pour
\regle{stale-closure}, son premier paragraphe reprend ce que nous avons
constaté, et la correction proposée est celle que nous avons appliquée :

\begin{sortie}
stale-closure
  long-lived callback keeps a state value frozen at registration time

A callback handed to `setInterval`, `addEventListener`, `subscribe`, `setTimeout` or a promise `.then` inside an effect closes over the variable values from the render that last ran the effect. [...]

Example:
  const [n, setN] = useState(0);
  useEffect(() => { setInterval(() => setN(n + 1), 1000); }, []);

Fix:
  Use the functional updater (`setN(n => n + 1)`) so the callback never reads the captured value; mirror the latest value into a `useRef` and read `ref.current`; or add the value to the deps array and return a cleanup so the callback is re-registered with a fresh capture.
\end{sortie}
(Sortie élaguée : la suite de l'explication, marquée \code{[...]}, détaille
la condition de l'\Error.)

\section{Le projet en chiffres}
\label{sec:introduction-chiffres}

Avant d'entrer dans le code, prenons la mesure de ce qu'il y a à lire. Tous les
chiffres de cette section sont comptés sur le dépôt au commit
\snapshotCommit{} (\code{find}, \code{wc -l}, \code{git rev-list}).

\subsection{Le code}

La bibliothèque \code{reactant} (\file{src/}) compte 142 fichiers Rust et
59\,757 lignes, commentaires et tests unitaires compris.
\Cref{tab:introduction-modules} les répartit par module. Deux modules pèsent à
eux seuls plus de la moitié du total : \file{src/rules} (les règles, leur API
typée, le front-end des règles déclaratives) et \file{src/engine} (le point
fixe, les relations, l'analyse de programme).

\begin{table}[htbp]\centering\small
\begin{tabularx}{\linewidth}{@{}lrrX@{}}
  \toprule
  Module & Fichiers & Lignes & Rôle (chapitres) \\
  \midrule
  \file{src/lowering} & 13 & 6\,692 & détection et lowering AST → IR (\ref{chap:frontend-detection}, \ref{chap:lowering-cfg}) \\
  \file{src/ir}       & 16 & 4\,423 & représentation intermédiaire, splice (\ref{chap:ir}, \ref{chap:splice-inlining}) \\
  \file{src/domains}  & 21 & 7\,393 & domaines abstraits, transfert, stores (\ref{chap:treillis-simples}--\ref{chap:transfert-stores}) \\
  \file{src/engine}   & 22 & 13\,913 & point fixe, relations, churn, programme (\ref{chap:cfg-analyzer-dominance}--\ref{chap:inter-composants}) \\
  \file{src/rules}    & 41 & 20\,432 & API des règles, règles natives, packs (\ref{chap:api-regles}--\ref{chap:regles-declaratives}) \\
  \file{src/resolver} & 4 & 1\,930 & découverte, parse, résolution d'imports (\ref{chap:projet-cli-driver}) \\
  \file{src/project}  & 4 & 1\,432 & projets Vite, Next.js, tsconfig (\ref{chap:projet-cli-driver}) \\
  \file{src/driver}   & 7 & 1\,666 & \code{run_check}, rendus humain et JSON (\ref{chap:projet-cli-driver}) \\
  \file{src/registry} & 3 & 691 & résumés de hooks de bibliothèques (\ref{chap:splice-inlining}) \\
  \file{src/cli}      & 7 & 570 & sous-commandes, chargement de la configuration (\ref{chap:projet-cli-driver}) \\
  racine de \file{src/} & 4 & 615 & \file{config.rs}, \file{lib.rs}, \file{main.rs}, \file{test_support.rs} \\
  \midrule
  total \file{src/}   & 142 & 59\,757 & \\
  \bottomrule
\end{tabularx}
\caption{Répartition du code de \file{src/} par module au commit
\snapshotCommit{} (fichiers \code{.rs}, lignes brutes).}
\label{tab:introduction-modules}
\end{table}

Autour de la bibliothèque :
\begin{itemize}
  \item \file{tests/} : 69 fichiers Rust de tests d'intégration
    (26\,212 lignes) et 131 fixtures (programmes React d'exemple, projets
    miniatures). Au total, le dépôt déclare 1\,529 fonctions \code{\#[test]}
    (632 dans \file{src/}, 897 dans \file{tests/}).
  \item \file{crates/reactant-wasm} : le point d'entrée WebAssembly
    (323 lignes), et \file{npm/} : le paquet \code{reactant-analyzer}
    (\cref{chap:distribution}).
  \item \file{docs/} : l'usage (\file{usage.md}), la page des limites
    (\file{limitations.md}), les règles personnalisées
    (\file{custom-rules.md}), le journal de précision, et surtout
    \file{docs/adr/} : \textbf{42 ADR} (\emph{Architecture Decision
    Records}), numérotées de \code{ADR-001} (\og concrete semantics\fg{}) à
    \code{ADR-042} (\og relations are engine products\fg{}). Chacune fait l'objet
    d'une fiche en annexe (\cref{chap:fiches-adr}).
\end{itemize}

L'historique compte 322 commits, du 8 mai au 27 septembre 2026, et six
versions publiées (\code{v0.1.0} à \code{v0.6.0}, la version de
\file{Cargo.toml}). Le \cref{chap:histoire-doctrine} raconte cette histoire ;
retenons-en qu'elle commence par une réécriture complète, décidée par la
conception (\adr{3} : une IR dédiée plutôt que des fonctions de transfert sur
l'AST), et que le pipeline actuel a été posé en une journée, le 1\textsuperscript{er}
juin, dans l'ordre même où ce manuscrit le présente.

\subsection{Les règles}

\reactant{} enregistre 19 règles natives (la fonction \code{all_rules},
\srcline{src/rules/mod.rs}{109}), qui émettent 21 \emph{noms de diagnostic} :
\regle{infinite-loop} émet aussi \regle{cross-component-infinite-loop}, et
\regle{setter-in-render} émet aussi \regle{cross-setter-in-render}.
\Cref{tab:introduction-catalogue} les liste tous, avec le résumé que donne
\code{reactant rules} et le niveau le plus haut que chaque nom peut atteindre
au commit \snapshotCommit{} (lu dans les constructeurs de diagnostic de
chaque règle, \file{src/rules/impls/}). Une configuration peut abaisser ces
niveaux, jamais les relever. L'annexe \ref{chap:catalogue-regles} donne une
fiche par règle.

\begin{table}[htbp]\centering\footnotesize
\begin{tabularx}{\linewidth}{@{}llX@{}}
  \toprule
  Diagnostic & Max & Résumé (\code{reactant rules}, verbatim) \\
  \midrule
  \regle{always-unstable-deps} & \Warning & a dep is a fresh reference every render, so the deps array never matches \\
  \regle{analysis-limit} & \Info & the analyzer deliberately truncated analysis here (potential false negatives) \\
  \regle{conditional-hook} & \Error & hook called inside a conditional branch \\
  \regle{cross-component-infinite-loop} & \Warning & child effect sets parent state, parent re-renders child, effect refires \\
  \regle{cross-setter-in-render} & \Error & parent's setter (received as prop) called during render \\
  \regle{derived-state} & \Warning & effect only mirrors another state, so compute it during render \\
  \regle{frozen-initial-state} & \Error & useState seeded from a prop that changes, so the state freezes at the first value \\
  \regle{infinite-loop} & \Error & effect sets state that re-triggers the effect, and the state diverges \\
  \regle{lazy-init} & \Error & useState initializer calls a function on every render \\
  \regle{missing-cleanup} & \Warning & effect starts something long-lived and returns no teardown \\
  \regle{missing-deps} & \Warning & effect body captures a variable not listed in its deps array \\
  \regle{redundant-set-state} & \Warning & setState called with the value the state already holds \\
  \regle{server-component-hook} & \Warning & a hook is called in a Next.js Server Component \\
  \regle{setter-in-render} & \Error & setState called during the render body \\
  \regle{stale-closure} & \Error & long-lived callback keeps a state value frozen at registration time \\
  \regle{state-lifted-too-high} & \Warning & state is used only deep in one child subtree; the components above re-render to pass it down \\
  \regle{state-mutation} & \Error & state or prop object mutated in place, same reference, no re-render \\
  \regle{unnecessary-rerender} & \Warning & mount-only effect immediately overwrites the initial state \\
  \regle{unstable-context-value} & \Warning & context provider hands consumers a new object every render \\
  \regle{wasted-subtree-render} & \Warning & a state written on every keystroke or pointer move re-renders subtrees that do not depend on it \\
  \regle{widening-info} & \Info & state slot required widening to converge (precision lost here) \\
  \bottomrule
\end{tabularx}
\caption{Le catalogue des diagnostics natifs au commit \snapshotCommit{}. La
colonne \og Max\fg{} est le niveau le plus haut atteignable ; pour
\regle{cross-component-infinite-loop}, le plafond \Warning{} est expliqué dans
le code (\src{src/rules/impls/infinite_loop.rs}{281}{283}).}
\label{tab:introduction-catalogue}
\end{table}

\begin{remarque}
  Un niveau maximal \Error{} ne veut pas dire que la règle émet souvent des
  \Error{} : c'est seulement le cas où une primitive de preuve a fourni son
  jeton. \regle{infinite-loop}, par exemple, émet un \Warning{} sur
  \code{Counter} et une \Error{} sur \code{Filters}.
\end{remarque}

À ces règles natives s'ajoutent des \emph{règles déclaratives} : des règles
écrites en JSON contre les relations du moteur, sans une ligne de Rust,
regroupées en \emph{packs} (\file{packs/guardrails.json} et cinq packs sous
\file{packs/community/}). Le test \file{tests/catalogue.rs} mesure
l'expressivité de ce langage sur un catalogue figé de 22 classes de règles :
21 sont exprimables au commit \snapshotCommit{}
(\code{EXPRESSIBLE_NOW = 21}, \srcline{tests/catalogue.rs}{1094} ;
\cref{chap:regles-declaratives}).

\subsection{Le corpus}

\begin{sloppypar}
Les tests unitaires vérifient des mécanismes ; ils ne disent pas ce que
l'outil produit sur du vrai code. Pour cela, le projet épingle un
\emph{corpus} de 14 dépôts open source (\code{ai-chatbot},
\code{bulletproof-react}, \code{chakra-ui}, \code{commerce}, \code{dub},
\code{excalidraw}, \code{mantine}, \code{memos},
\code{next-shadcn-dashboard-starter}, \code{novel}, \code{precedent},
\code{shadcn-admin}, \code{twenty}, \code{zustand}), clonés à des commits
fixés par \file{scripts/setup-test-repo.sh}. Le nombre de diagnostics
distincts produits sur ce corpus est commité dans
\file{docs/corpus-baseline.json} : \textbf{1\,498} localisations au commit
\snapshotCommit{}, dont 476 \regle{missing-deps}, 381
\regle{always-unstable-deps}, 222 \regle{lazy-init} et 35
\regle{infinite-loop}. Un job d'intégration continue compare chaque
exécution à ce chiffre ; toute variation doit être expliquée, et consignée
dans \file{docs/precision-log.md}. Le \cref{chap:tests-corpus} décrit cette
méthode, et pourquoi un chiffre qui \emph{baisse} n'est pas forcément une
bonne nouvelle : pour un outil sound, un diagnostic qui disparaît peut être un
faux positif corrigé, ou un faux négatif introduit.
\end{sloppypar}

\section{Plan du manuscrit et conseils de lecture}
\label{sec:introduction-plan}

Le manuscrit suit le trajet d'un fichier source dans le pipeline de la
\cref{fig:introduction-pipeline}, précédé du contexte nécessaire et suivi de
l'outil et de sa méthode. Il est organisé en six parties et quatre annexes
(\cref{tab:introduction-parties}) ; la difficulté croît d'une partie à
l'autre, et à l'intérieur de chaque chapitre.

\begin{table}[htbp]\centering\small
\begin{tabularx}{\linewidth}{@{}lXc@{}}
  \toprule
  Partie & Contenu & Difficulté \\
  \midrule
  I. Le problème et son contexte & ce chapitre ; le modèle d'exécution de React (\ref{chap:react-modele}) ; rappels d'interprétation abstraite (\ref{chap:interpretation-abstraite}) ; tour d'architecture (\ref{chap:tour-architecture}) & 1--2 \\
  II. Du source à l'IR & oxc et détection (\ref{chap:frontend-detection}), IR (\ref{chap:ir}), lowering (\ref{chap:lowering-cfg}), splice et inlining (\ref{chap:splice-inlining}) & 2--3 \\
  III. Les domaines abstraits & treillis simples (\ref{chap:treillis-simples}), produit \code{StateValue} et stabilité (\ref{chap:statevalue-stabilite}), transfert et stores (\ref{chap:transfert-stores}) & 2--4 \\
  IV. Le moteur & CFG et dominance (\ref{chap:cfg-analyzer-dominance}), point fixe d'un composant (\ref{chap:point-fixe-composant}), relations (\ref{chap:relations}), churn (\ref{chap:churn}), inter-composants (\ref{chap:inter-composants}) & 3--5 \\
  V. Les règles & API typée (\ref{chap:api-regles}), règles d'état (\ref{chap:regles-etat}), \regle{infinite-loop} (\ref{chap:infinite-loop}), effets (\ref{chap:regles-effets}), rendu et RSC (\ref{chap:regles-rendu-rsc}), règles déclaratives (\ref{chap:regles-declaratives}) & 3--5 \\
  VI. L'outil et sa méthode & CLI et driver (\ref{chap:projet-cli-driver}), distribution (\ref{chap:distribution}), tests et corpus (\ref{chap:tests-corpus}), histoire et doctrine (\ref{chap:histoire-doctrine}), limites (\ref{chap:limites-perspectives}) & 2--3 \\
  Annexes & fiches ADR (\ref{chap:fiches-adr}), glossaire (\ref{chap:glossaire}), catalogue des règles (\ref{chap:catalogue-regles}) & --- \\
  \bottomrule
\end{tabularx}
\caption{Les parties du manuscrit ; les numéros entre parenthèses renvoient aux chapitres.}
\label{tab:introduction-parties}
\end{table}

\subsection{Parcours de lecture}

Le manuscrit se lit dans l'ordre, mais tous les lecteurs n'ont pas les mêmes
besoins. Quelques parcours :

\begin{description}
  \item[Vous utilisez l'outil et voulez comprendre ses diagnostics.] Ce
    chapitre, puis le \cref{chap:react-modele} ; ensuite, pour chaque
    diagnostic rencontré, sa fiche dans l'annexe \ref{chap:catalogue-regles}
    et sa section dans les \cref{chap:regles-etat,chap:infinite-loop,chap:regles-effets,chap:regles-rendu-rsc} ;
    enfin le \cref{chap:projet-cli-driver} (configuration, projets) et le
    \cref{chap:limites-perspectives} (ce que l'outil rate).
  \item[Vous voulez écrire une règle.] La partie~I, puis le
    \cref{chap:api-regles} (l'API typée et le jeton de preuve), le
    \cref{chap:relations} (les faits disponibles), une ou deux règles de la
    partie~V comme modèles, et le \cref{chap:regles-declaratives} si une règle
    déclarative suffit. Le \cref{chap:tests-corpus} dit comment mesurer
    l'effet d'une nouvelle règle.
  \item[Vous voulez modifier le moteur.] Tout, dans l'ordre, avec une
    attention particulière aux \cref{chap:point-fixe-composant,chap:relations,chap:churn}
    et aux encadrés \og piège\fg{} qui signalent les trous de soundness
    connus.
  \item[Vous connaissez l'analyse statique, pas React.] Le
    \cref{chap:react-modele} est indispensable ; le
    \cref{chap:interpretation-abstraite} peut être survolé pour la notation.
    Les chapitres propres au projet sont ensuite les
    \cref{chap:statevalue-stabilite,chap:point-fixe-composant,chap:churn,chap:inter-composants}.
  \item[Vous connaissez React, pas l'analyse statique.] Le
    \cref{chap:interpretation-abstraite} est indispensable : il pose le
    vocabulaire (treillis, join, widening, must/may, soundness) que tout le
    reste emploie.
  \item[Vous vous intéressez aux choix de conception.] Le
    \cref{chap:histoire-doctrine} et les fiches ADR (\cref{chap:fiches-adr}),
    puis les encadrés \og décision\fg{} des chapitres.
\end{description}

\subsection{Conventions}

\begin{itemize}
  \item \textbf{Code cité.} Les extraits pris dans le dépôt portent en légende
    leur fichier et leurs vrais numéros de ligne au commit
    \snapshotCommit{} (par exemple \cref{lst:introduction-severity}). Les
    programmes React d'exemple sont écrits pour le manuscrit ; leur chemin
    sous \file{/tmp} est indiqué pour pouvoir les réanalyser. Un extrait
    élagué le dit dans sa légende.
  \item \textbf{Sorties.} Toute sortie de l'analyseur a été observée
    réellement sur ce commit ; les élisions sont signalées. Elles sont
    reproduites en anglais, telles que l'outil les écrit.
  \item \textbf{Encadrés.} \og React\fg{} rappelle un comportement de React
    nécessaire à la suite ; \og Décision\fg{} résume une ADR et les
    alternatives écartées ; \og Piège\fg{} signale une erreur naturelle ou une
    subtilité, y compris les écarts observés entre la visée et le code.
  \item \textbf{Vocabulaire.} Les noms du code restent en anglais
    (\code{StateValue}, \code{SlotWriter}) ; les concepts ont un nom
    français fixé, recensé dans le glossaire (\cref{chap:glossaire}). Quelques
    mots anglais sont gardés tels quels parce qu'ils sont d'usage :
    \emph{widening}, \emph{join}, \emph{meet}, \emph{soundness}, \emph{slot},
    \emph{handler}, \emph{deps}.
\end{itemize}

\begin{resume}
  \begin{itemize}
    \item Les défauts React les plus coûteux --- boucle de rendus, dépendance
      oubliée, fermeture figée --- ne se voient pas au premier rendu : ils
      sont des propriétés de la \emph{suite} des rendus, commits et effets
      (\cref{tab:introduction-trois-bugs}).
    \item Un linter syntaxique juge la forme d'une fonction, dans un fichier.
      Il voit un hook conditionnel ou un dep absent ; il ne voit ni la suite
      des valeurs d'un état, ni un alias de setter, ni un hook personnalisé
      dans son appelant, ni un setter passé en prop
      (\cref{tab:introduction-syntaxe}).
    \item \reactant{} interprète abstraitement chaque composant : il rejoue
      les tours de React sur des \emph{ensembles} de valeurs (intervalles,
      stabilité des références, écrivains de chaque état) jusqu'à un point
      fixe, que le widening garantit fini ; les règles lisent ce résultat.
    \item Pipeline : parse (\code{oxc}) → lowering → IR (\code{ComponentIR})
      → moteur (\code{AnalysisResult}) → relations → règles
      (\code{Diagnostic}) → driver (sortie humaine ou JSON)
      (\cref{fig:introduction-pipeline}).
    \item Trois niveaux : \Error{} = défaut certain, prouvé en entier (seul
      \code{Diagnostic::error}, qui exige un \code{Certified}, en produit) ;
      \Warning{} = défaut possible ou coût incertain ; \Info{} = limite de
      l'analyse. Sous-déclarer est permis (\code{Counter}, \issue{144}),
      sur-déclarer ne l'est pas.
    \item Soundness : l'analyse sur-approxime ; FP tolérés, FN interdits. Le
      silence est une affirmation (\code{verified}), et l'outil refuse le
      \og no issues found\fg{} sur du code qu'il n'a pas lu. C'est une visée
      du code, avec des écarts connus et documentés
      (\cref{lst:introduction-fn-updater}).
    \item \sloppy Fichiers rencontrés :
      \file{src/rules/api/diagnostic.rs} (\code{Severity},
      \code{Diagnostic::error}) ; \file{src/rules/mod.rs}
      (\code{all_rules}) ; \file{src/driver/mod.rs} (\code{run_check}) ;
      \file{docs/usage.md} ; \file{docs/adr/}.
    \item Chiffres au commit \snapshotCommit{} : 142 fichiers et
      59\,757 lignes dans \file{src/}, 1\,529 tests, 19 règles natives et
      21 noms de diagnostic, 42 ADR, 14 dépôts de corpus, 1\,498
      localisations de référence.
  \end{itemize}
  Le \cref{chap:react-modele} précise le modèle d'exécution de React que
  l'analyse abstrait ; le \cref{chap:interpretation-abstraite} donne les
  définitions exactes de ce que ce chapitre a présenté en mots.
\end{resume}

\section{Exercices}
\label{sec:introduction-exercices}

Les sorties citées dans les solutions ont été observées avec
\code{check --no-color --info --trace} sur les fichiers indiqués.

\begin{exercice}{Un seuil plus lointain}{introduction-hundred}
  On remplace la garde de \code{Bounded} (\cref{lst:introduction-bounded})
  par \code{if (n < 100) setN(n + 1)} (\file{/tmp/ch01/exo_hundred.tsx}).
  (a) Que fait React ? (b) En raisonnant comme sur la
  \cref{fig:introduction-intervalles}, quel intervalle final attendez-vous
  pour \code{n} ? (c) \regle{infinite-loop} doit-elle se déclencher ?
\end{exercice}
\begin{solution}
  (a) Cent rendus supplémentaires après le montage, puis l'effet cesse
  d'écrire : \code{n} s'arrête à \code{100}. (b) Le widening projette la
  borne qui grandit sur le plus petit seuil qui l'englobe ; les littéraux du
  composant fournissent \code{100}, d'où $[0,100]$, stable parce que l'effet
  n'écrit que sous $n<100$. Le moteur calcule bien \code{number[0, 100]}
  (sonde hors dépôt). (c) Non : la sortie contient l'\Info{}
  \regle{widening-info} (l'état a été élargi) et la ligne \code{verified
  infinite-loop no effect diverges into an infinite render loop}. Le
  widening seul n'est pas un diagnostic de boucle.
\end{solution}

\begin{exercice}{Élargi, et pourtant sans boucle}{introduction-once}
  Soit l'effet \code{useEffect(() => \{ setN(n + 1); \}, [])} dans un
  composant \code{Once} dont l'état vaut initialement \code{0}
  (\file{/tmp/ch01/exo_once.tsx}). (a) Combien de rendus React fait-il ?
  (b) L'analyseur élargit l'état \code{n} à $[0,+\infty)$ : pourquoi, alors
  que l'effet ne s'exécute qu'une fois ? (c) Quels diagnostics attendez-vous ?
\end{exercice}
\begin{solution}
  (a) Deux : le montage (\code{n = 0}), puis le rendu provoqué par
  \code{setN(1)} ; les deps \code{[]} n'ayant pas changé, l'effet ne repart
  pas. (b) Le tour abstrait exécute \emph{tous} les effets à chaque tour,
  qu'ils soient re-déclenchés ou non (\cref{sec:introduction-pari}) : c'est
  une sur-approximation, qui fait croître \code{n} à chaque tour abstrait. La
  règle \regle{infinite-loop}, elle, tient compte des deps et ne conclut pas à
  une boucle. (c) La sortie observée contient un \Warning{}
  \regle{missing-deps} sur \code{n} (lu, non déclaré), l'\Info{}
  \regle{widening-info}, et \code{verified infinite-loop}. Le code de sortie
  est \code{1}, à cause du \Warning.
\end{solution}

\begin{exercice}{Setter dans le corps du rendu}{introduction-eager}
  Quel niveau \reactant{} doit-il donner au composant suivant
  (\file{/tmp/ch01/exo_render.tsx}, ligne d'import omise), et pourquoi ?
\begin{lstlisting}[language=TypeScript]
export function Eager() {
  const [n, setN] = useState(0);
  setN(n + 1);
  return <p>{n}</p>;
}
\end{lstlisting}
\end{exercice}
\begin{solution}
  Un setter appelé pendant le rendu planifie un nouveau rendu avant même le
  commit ; appelé inconditionnellement, il relance le rendu à chaque fois, et
  React interrompt la boucle par une erreur (\cref{chap:react-modele}). Le
  défaut est certain dès que le composant est rendu, et la preuve est
  structurelle : l'appel est dans le bloc d'entrée du rendu, qui est sur tous
  les chemins. \reactant{} donne \code{error setter-in-render [hook:0]
  (line 5:2) setter `setN` called directly in the render body, move this
  call into a useEffect or an event handler} (\cref{chap:regles-etat}).
\end{solution}

\begin{exercice}{Une constante de module}{introduction-const}
  Dans \file{/tmp/ch01/exo_const.tsx}, un effet de deps \code{[page]} lit une
  prop \code{page} et une constante de module \code{const LIMIT = 10}.
  \regle{missing-deps} doit-elle signaler \code{LIMIT} ? Que répond l'outil ?
\end{exercice}
\begin{solution}
  Non : une constante de module ne change jamais entre deux rendus, l'omettre
  des deps ne peut pas figer une valeur périmée. L'analyse le sait parce
  qu'elle évalue les constantes de module (\cref{chap:frontend-detection}).
  Sortie observée : \code{1 clean component(s) hidden} puis
  ✓~\code{1 file(s) no issues found.}
\end{solution}

\begin{exercice}{Classer des énoncés}{introduction-classer}
  Pour chacun des énoncés suivants, dites s'il peut justifier une \Error, un
  \Warning{} ou seulement une \Info{} : (a) \og cet effet n'a pas de
  tableau de deps et écrit, sur tous ses chemins, un objet littéral neuf dans
  un état\fg{} ; (b) \og cet état a dû être élargi pendant l'analyse\fg{} ;
  (c) \og cette variable, lue dans l'effet, peut changer entre deux
  rendus\fg{} ; (d) \og ce hook importé n'a pas pu être lu\fg{}.
\end{exercice}
\begin{solution}
  (a) \Error{} possible : chaque conjoint est un fait certain (écriture sur
  tous les chemins, référence toujours neuve, re-déclenchement à chaque
  commit sans deps), la conclusion est prouvée en entier. Sur la variante de
  \code{Filters} (\cref{ex:introduction-churn}) sans tableau de deps
  (\file{/tmp/ch01/exo_nodeps_churn.tsx}), l'outil répond \code{error
  infinite-loop [hook:1] (line 5:2) this effect has no dependency array and
  stores a fresh reference into state `query`, so it re-runs after every
  render and re-triggers itself: infinite render loop}. (b) \Info{} seulement : un
  widening peut venir d'une imprécision, c'est l'\Info{}
  \regle{widening-info}. (c) \Warning{} : défaut possible, qui dépend de
  l'évolution réelle de la variable (\regle{missing-deps}). (d) \Info{} :
  c'est une limite de l'analyse (\regle{analysis-limit}), qui suspend aussi
  les lignes \code{verified} du composant.
\end{solution}

